% Introduction to the EU - Week 9 Reading Summary
% Compile: pdflatex European-Parliament-reading-summary.tex
%
% Fill-in rules:
% - Replace placeholders with exact wording from the PDF (do not guess).
% - For workshop relevance, focus on "housing crisis" frames/lexicon (as they appear in the text).

\documentclass[11pt,a4paper]{article}
\usepackage[utf8]{inputenc}
\usepackage[T1]{fontenc}
\usepackage[margin=2.5cm]{geometry}
\usepackage{booktabs}
\usepackage{enumitem}
\setlist{nosep,leftmargin=*}
\usepackage{parskip}
\usepackage{microtype}
\usepackage{hyperref}
\usepackage{xcolor}
\definecolor{eublue}{RGB}{0,51,153}

\hyphenation{Spitzenkandidat Spitzenkandidaten}
\hypersetup{
  colorlinks=true,
  linkcolor=eublue,
  urlcolor=eublue,
  pdftitle={Reading Summary},
  pdfauthor={Introduction to the EU, UCM}
}

\begin{document}

\title{Author Name: \textit{European Parliament}\\Key Themes \& Arguments --- Housing-crisis workshop lens}
\author{Introduction to the European Union\\Universidad Complutense de Madrid\\Week 9}
\date{}
\maketitle
\vspace{-1em}
\hrule
\vspace{1em}

\section*{Rhetorical Framework (Ethos, Logos, Pathos)}
\begin{table}[h]
\centering
\begin{tabular}{@{}lll@{}}
\toprule
\textbf{Appeal} & \textbf{Meaning} & \textbf{Author's Use} \\
\midrule
\textbf{Ethos} & Credibility, authority, trust & \\
\textbf{Logos} & Logic, structure, argument & \\
\textbf{Pathos} & Emotion, stakes, persuasion & \\
\bottomrule
\end{tabular}
\end{table}

\textbf{Key insight:} One paragraph on how the author blends ethos, logos, pathos (fill with your extracted points).

\section*{Bibliographic Information}
\begin{itemize}
\item \textbf{Author:} [paste]
\item \textbf{Title:} European Parliament (as shown on PDF)
\item \textbf{Source:} European Parliament.pdf
\item \textbf{Year:} [paste]
\end{itemize}

\section*{Summary & Key Points}
\subsection*{I. First Theme}
Content...

\subsection*{II. Second Theme}
Content...

\section*{Main Arguments}
\subsection*{I. First Argument}
Content...

\subsection*{II. Second Argument}
Content...

\section*{Context & Analysis}
\begin{itemize}
\item \textbf{Historical:} [what context the text references]
\item \textbf{Connections:} [connect to workshop lens: housing crisis framing, if present]
\end{itemize}

\section*{Relevance to Course}
\begin{itemize}
\item \textbf{Ethos:} [why it matters / credibility]
\item \textbf{Logos:} [how it informs reasoning about institutions/procedure]
\item \textbf{Pathos:} [what stakes it emphasizes]
\item \textbf{All three:} [one-line integrated relevance]
\end{itemize}

\section*{Discussion Questions}
\begin{enumerate}
\item Which terms/lexicon does the text use to define "the housing problem" (if housing is discussed), and what framing effect does that create?
\item How does the text describe parliamentary procedure/roles, and how would that shape a housing-crisis debate?
\item Which group-oriented argument patterns (lexicon/tools/EU role stance) would this text help you justify during the workshop?
\end{enumerate}

\vspace{1em}
\noindent\footnotesize\textit{Introduction to the European Union, Universidad Complutense de Madrid.}

\end{document}

